Imagine you are an observer agent living in a universe with a smooth continuum of visible angular directions. Your visible sky is modeled as the complex projective line $\mathbb P^1(\mathbb C)$, equipped with normalized Fubini-Study probability measure $\sigma$. You as an observer are not told whether this continuum is fundamental, or whether it is the low-energy rendering/simulation of a hidden finite algebraic substrate. If you are a capable observer to figure out whether you live in a simulation or not, you should be able to infer, from a small statistical defect in the visible angular distribution together with a few invariant-theoretic fingerprints, whether the continuum hypothesis is best explained by an underlying quotient substrate. You not given the hidden group, the primitive invariant, or the perturbation strength. These must be reconstructed from the algebraic data and then used in an asymptotic decision problem.

Let $U=\mathbb C^2$, with homogeneous coordinates $(u,v)$, and let $\sigma$ be the normalized Fubini-Study probability measure on \(\mathbb P^1(\mathbb C)\). Under the alternative hypothesis, a hidden computational substrate is assumed to have a non-cyclic finite \(\operatorname{SL}_2(\mathbb C)\)-quotient singularity $X_\Gamma=\operatorname{Spec} \mathbb C[u,v]^\Gamma$, for some unknown finite subgroup $\Gamma\subset \operatorname{SL}_2(\mathbb C).$ The visible angular sky is \(\mathbb P^1(\mathbb C)\), and the local statistical defect is generated by the lowest primitive invariant of the quotient singularity.

The substrate is known only through the following side information.

1. The minimal resolution of $X_\Gamma$ has exceptional intersection graph $E_8$.

2. Equivalently, $\mathbb C[u,v]^\Gamma$ has primitive homogeneous generators of degrees $12,20,30$, satisfying one weighted-homogeneous relation of degree $60$.

3. The degree-$12$ primitive invariant is normalized as follows. It lies in the one-parameter family
\[
F_a(u,v)=u^{11}v+a\,u^6v^6-uv^{11},
\]
and \(a\in \mathbb C\) is chosen so that, if
\[
H_a=\operatorname{Hess}(F_a)
=
\det
\begin{pmatrix}
\dfrac{\partial^2 F_a}{\partial u^2}
&
\dfrac{\partial^2 F_a}{\partial u\partial v}
\\[2mm]
\dfrac{\partial^2 F_a}{\partial v\partial u}
&
\dfrac{\partial^2 F_a}{\partial v^2}
\end{pmatrix},
\]
and
\[
T_a=\operatorname{Jac}(F_a,H_a)
=
\det
\begin{pmatrix}
\dfrac{\partial F_a}{\partial u}
&
\dfrac{\partial F_a}{\partial v}
\\[2mm]
\dfrac{\partial H_a}{\partial u}
&
\dfrac{\partial H_a}{\partial v}
\end{pmatrix},
\]
then the three degree-\(60\) binary forms $F_a^5, H_a^3, T_a^2$, are linearly dependent over $\mathbb C$, with all three coefficients in the linear dependence nonzero. This is the intrinsic syzygy condition seen by the observer.

For such an admissible value of $a$, define the projectively well-defined nonnegative function
\[
Q([u:v])=
\frac{|F_a(u,v)|^2}{\bigl(|u|^2+|v|^2\bigr)^{12}},
\qquad [u:v]\in \mathbb P^1(\mathbb C),
\]
and its mean-zero normalization
\[
\rho([u:v])=
\frac{Q([u:v])}{\displaystyle\int_{\mathbb P^1} Q\,d\sigma}-1.
\]
For sufficiently small \(\varepsilon\), define a probability measure \(P_\varepsilon\) on \(\mathbb P^1(\mathbb C)\) by
\[
\frac{dP_\varepsilon}{d\sigma}=1+\varepsilon\rho.
\]

The null hypothesis is that the visible sky is statistically indistinguishable from the continuum model, $H_0:\varepsilon=0.$ The alternative hypothesis is that the observed sky is the continuum rendering of the hidden quotient substrate, $H_1:\varepsilon=\varepsilon_0>0.$
You as an observer see $n$ independent samples from $P_\varepsilon$, where $\varepsilon\in\{0,\varepsilon_0\}.$ You as an observer scientist use the Neyman/Pearson likelihood-ratio test for
\[
H_0:\varepsilon=0
\qquad\text{against}\qquad
H_1:\varepsilon=\varepsilon_0>0.
\]
In the local regime $\varepsilon_0=\frac{h}{\sqrt n},$ determine the exact constant \(C\) such that, when $h=C,$ the asymptotic Type-I and Type-II errors are both $\alpha=\beta=\Phi(-1).$

Your answer must include the following five items:

1. the hidden group $\Gamma$, up to conjugacy in $\operatorname{SL}_2(\mathbb C)$

2.  the admissible value or values of $a$

3. the exact Fisher information $\mathcal I=\int_{\mathbb P^1}\rho^2\,d\sigma$

4. the exact threshold constant $C=\frac{2}{\sqrt{\mathcal I}}$

5. the limiting log-Bayes factor, in favor of the hidden quotient-substrate sky over the continuum sky, assuming equal prior odds.
